\section{System Design and Implementation}

This section details the layered architecture, the implementation of each pipeline stage,
the AI routing layer, and the configuration system.

% ============================================================
\subsection{Layered Architecture}
% ============================================================

Geo-SLM follows a strict 5-layer separation, as shown in \Cref{fig:layer_arch}. The
core engine (Layer~3) has \textbf{zero dependency} on any web framework; it runs
identically as a standalone Python library, a CLI tool, or behind a FastAPI server.

\begin{figure}[htbp]
\centering
% TikZ: System Layer Architecture
% Usage: % TikZ: System Layer Architecture
% Usage: % TikZ: System Layer Architecture
% Usage: \input{figures/tikz/tikz_layer_architecture.tex}
% Requires: \usepackage{tikz} \usetikzlibrary{shapes.geometric, arrows.meta, positioning, fit, backgrounds, calc}

\begin{tikzpicture}[
    node distance=0.5cm,
    layer/.style={
        rectangle, rounded corners=2pt, draw=#1!60!black, thick,
        minimum width=13cm, minimum height=1.4cm,
        fill=#1!8, font=\small\bfseries, align=center
    },
    component/.style={
        rectangle, rounded corners=2pt, draw=black!40,
        minimum width=2.8cm, minimum height=0.7cm,
        fill=white, font=\footnotesize, align=center
    },
    arrow/.style={-{Stealth[length=3mm]}, thick, draw=black!50},
    layerlabel/.style={font=\footnotesize\bfseries, text=#1!70!black}
]

% Layer 1: Interface
\node[layer=blue] (l1) {Interface Layer};
\node[component, above=0.15cm of l1.south west, anchor=south west, xshift=0.5cm] (cli) {CLI (Typer)};
\node[component, right=0.3cm of cli] (demo) {Streamlit Demo};
\node[component, right=0.3cm of demo] (api) {FastAPI Server};
\node[component, right=0.3cm of api] (notebook) {Jupyter Notebook};

% Layer 2: Task Queue
\node[layer=teal, below=0.6cm of l1] (l2) {Task Queue Layer};
\node[component, above=0.15cm of l2.south west, anchor=south west, xshift=2.0cm] (celery) {Celery Workers};
\node[component, right=0.4cm of celery] (redis) {Redis Broker};
\node[component, right=0.4cm of redis] (pg) {PostgreSQL State};

% Layer 3: Core Engine
\node[layer=orange, below=0.6cm of l2] (l3) {Core Engine (Pure Python)};
\node[component, above=0.15cm of l3.south west, anchor=south west, xshift=0.3cm] (pipeline) {Pipeline\\Orchestrator};
\node[component, right=0.2cm of pipeline] (stages) {Stages\\S1--S5};
\node[component, right=0.2cm of stages] (schemas) {Pydantic\\Schemas};
\node[component, right=0.2cm of schemas] (router) {AI Router};
\node[component, right=0.2cm of router] (validators) {Validators};

% Layer 4: Data
\node[layer=purple, below=0.6cm of l3] (l4) {Data \& Model Layer};
\node[component, above=0.15cm of l4.south west, anchor=south west, xshift=0.5cm] (files) {File Storage};
\node[component, right=0.3cm of files] (cache) {OCR Cache};
\node[component, right=0.3cm of cache] (weights) {Model Weights};
\node[component, right=0.3cm of weights] (config) {YAML Config};

% Arrows
\draw[arrow] (l1) -- (l2);
\draw[arrow] (l2) -- (l3);
\draw[arrow] (l3) -- (l4);

% Side labels
\node[layerlabel=blue, rotate=90, anchor=south] at ($(l1.west) + (-0.5, 0)$) {Users};
\node[layerlabel=teal, rotate=90, anchor=south] at ($(l2.west) + (-0.5, 0)$) {Async};
\node[layerlabel=orange, rotate=90, anchor=south] at ($(l3.west) + (-0.5, 0)$) {Logic};
\node[layerlabel=purple, rotate=90, anchor=south] at ($(l4.west) + (-0.5, 0)$) {Storage};

\end{tikzpicture}

% Requires: \usepackage{tikz} \usetikzlibrary{shapes.geometric, arrows.meta, positioning, fit, backgrounds, calc}

\begin{tikzpicture}[
    node distance=0.5cm,
    layer/.style={
        rectangle, rounded corners=2pt, draw=#1!60!black, thick,
        minimum width=13cm, minimum height=1.4cm,
        fill=#1!8, font=\small\bfseries, align=center
    },
    component/.style={
        rectangle, rounded corners=2pt, draw=black!40,
        minimum width=2.8cm, minimum height=0.7cm,
        fill=white, font=\footnotesize, align=center
    },
    arrow/.style={-{Stealth[length=3mm]}, thick, draw=black!50},
    layerlabel/.style={font=\footnotesize\bfseries, text=#1!70!black}
]

% Layer 1: Interface
\node[layer=blue] (l1) {Interface Layer};
\node[component, above=0.15cm of l1.south west, anchor=south west, xshift=0.5cm] (cli) {CLI (Typer)};
\node[component, right=0.3cm of cli] (demo) {Streamlit Demo};
\node[component, right=0.3cm of demo] (api) {FastAPI Server};
\node[component, right=0.3cm of api] (notebook) {Jupyter Notebook};

% Layer 2: Task Queue
\node[layer=teal, below=0.6cm of l1] (l2) {Task Queue Layer};
\node[component, above=0.15cm of l2.south west, anchor=south west, xshift=2.0cm] (celery) {Celery Workers};
\node[component, right=0.4cm of celery] (redis) {Redis Broker};
\node[component, right=0.4cm of redis] (pg) {PostgreSQL State};

% Layer 3: Core Engine
\node[layer=orange, below=0.6cm of l2] (l3) {Core Engine (Pure Python)};
\node[component, above=0.15cm of l3.south west, anchor=south west, xshift=0.3cm] (pipeline) {Pipeline\\Orchestrator};
\node[component, right=0.2cm of pipeline] (stages) {Stages\\S1--S5};
\node[component, right=0.2cm of stages] (schemas) {Pydantic\\Schemas};
\node[component, right=0.2cm of schemas] (router) {AI Router};
\node[component, right=0.2cm of router] (validators) {Validators};

% Layer 4: Data
\node[layer=purple, below=0.6cm of l3] (l4) {Data \& Model Layer};
\node[component, above=0.15cm of l4.south west, anchor=south west, xshift=0.5cm] (files) {File Storage};
\node[component, right=0.3cm of files] (cache) {OCR Cache};
\node[component, right=0.3cm of cache] (weights) {Model Weights};
\node[component, right=0.3cm of weights] (config) {YAML Config};

% Arrows
\draw[arrow] (l1) -- (l2);
\draw[arrow] (l2) -- (l3);
\draw[arrow] (l3) -- (l4);

% Side labels
\node[layerlabel=blue, rotate=90, anchor=south] at ($(l1.west) + (-0.5, 0)$) {Users};
\node[layerlabel=teal, rotate=90, anchor=south] at ($(l2.west) + (-0.5, 0)$) {Async};
\node[layerlabel=orange, rotate=90, anchor=south] at ($(l3.west) + (-0.5, 0)$) {Logic};
\node[layerlabel=purple, rotate=90, anchor=south] at ($(l4.west) + (-0.5, 0)$) {Storage};

\end{tikzpicture}

% Requires: \usepackage{tikz} \usetikzlibrary{shapes.geometric, arrows.meta, positioning, fit, backgrounds, calc}

\begin{tikzpicture}[
    node distance=0.5cm,
    layer/.style={
        rectangle, rounded corners=2pt, draw=#1!60!black, thick,
        minimum width=13cm, minimum height=1.4cm,
        fill=#1!8, font=\small\bfseries, align=center
    },
    component/.style={
        rectangle, rounded corners=2pt, draw=black!40,
        minimum width=2.8cm, minimum height=0.7cm,
        fill=white, font=\footnotesize, align=center
    },
    arrow/.style={-{Stealth[length=3mm]}, thick, draw=black!50},
    layerlabel/.style={font=\footnotesize\bfseries, text=#1!70!black}
]

% Layer 1: Interface
\node[layer=blue] (l1) {Interface Layer};
\node[component, above=0.15cm of l1.south west, anchor=south west, xshift=0.5cm] (cli) {CLI (Typer)};
\node[component, right=0.3cm of cli] (demo) {Streamlit Demo};
\node[component, right=0.3cm of demo] (api) {FastAPI Server};
\node[component, right=0.3cm of api] (notebook) {Jupyter Notebook};

% Layer 2: Task Queue
\node[layer=teal, below=0.6cm of l1] (l2) {Task Queue Layer};
\node[component, above=0.15cm of l2.south west, anchor=south west, xshift=2.0cm] (celery) {Celery Workers};
\node[component, right=0.4cm of celery] (redis) {Redis Broker};
\node[component, right=0.4cm of redis] (pg) {PostgreSQL State};

% Layer 3: Core Engine
\node[layer=orange, below=0.6cm of l2] (l3) {Core Engine (Pure Python)};
\node[component, above=0.15cm of l3.south west, anchor=south west, xshift=0.3cm] (pipeline) {Pipeline\\Orchestrator};
\node[component, right=0.2cm of pipeline] (stages) {Stages\\S1--S5};
\node[component, right=0.2cm of stages] (schemas) {Pydantic\\Schemas};
\node[component, right=0.2cm of schemas] (router) {AI Router};
\node[component, right=0.2cm of router] (validators) {Validators};

% Layer 4: Data
\node[layer=purple, below=0.6cm of l3] (l4) {Data \& Model Layer};
\node[component, above=0.15cm of l4.south west, anchor=south west, xshift=0.5cm] (files) {File Storage};
\node[component, right=0.3cm of files] (cache) {OCR Cache};
\node[component, right=0.3cm of cache] (weights) {Model Weights};
\node[component, right=0.3cm of weights] (config) {YAML Config};

% Arrows
\draw[arrow] (l1) -- (l2);
\draw[arrow] (l2) -- (l3);
\draw[arrow] (l3) -- (l4);

% Side labels
\node[layerlabel=blue, rotate=90, anchor=south] at ($(l1.west) + (-0.5, 0)$) {Users};
\node[layerlabel=teal, rotate=90, anchor=south] at ($(l2.west) + (-0.5, 0)$) {Async};
\node[layerlabel=orange, rotate=90, anchor=south] at ($(l3.west) + (-0.5, 0)$) {Logic};
\node[layerlabel=purple, rotate=90, anchor=south] at ($(l4.west) + (-0.5, 0)$) {Storage};

\end{tikzpicture}

\caption{Geo-SLM 5-layer system architecture. Each layer depends only on layers below it.
    The Core Engine is independent of the serving layer.}
\label{fig:layer_arch}
\end{figure}

\Cref{tab:tech} lists the technology stack for each component.

\begin{table}[htbp]
\centering
\caption{Technology stack overview.}
\label{tab:tech}
\begin{table}[htbp]
\centering
\caption{Core Technology Stack}
\label{tab:technology-stack}
\begin{tabular}{@{}llp{5cm}@{}}
\toprule
\textbf{Component} & \textbf{Technology} & \textbf{Rationale} \\
\midrule
Language         & Python 3.11+          & AI/ML ecosystem maturity \\
Object Detection & YOLOv8m (Ultralytics) & Fast, accurate, trainable \\
OCR              & EasyOCR / PaddleOCR   & Multi-language support \\
Image Processing & OpenCV + Pillow       & Industry standard \\
Geometric Calc.  & NumPy + Custom        & Precision arithmetic \\
Classification   & ResNet-18 (PyTorch)   & Lightweight, 94.14\% accuracy \\
Data Validation  & Pydantic v2           & Runtime schema enforcement \\
Configuration    & OmegaConf (YAML)      & Hierarchical ML config \\
\bottomrule
\end{tabular}
\end{table}
\end{table}

% ============================================================
\subsection{Stage 1: Ingestion and Sanitation}
% ============================================================

Stage~1 transforms diverse input formats (PDF, PNG, JPG, TIFF, BMP) into normalized
\texttt{CleanImage} objects. Key operations include:

\begin{itemize}
    \item \textbf{PDF rendering}: PyMuPDF renders each page using a transformation matrix
          \texttt{Matrix(DPI/72, DPI/72)} at a configurable DPI (default: 150).
    \item \textbf{Quality validation}: Blur detection using Laplacian variance
          ($\sigma^2(\nabla^2 I)$); images below the threshold (100.0) are flagged but
          not rejected, deferring the decision to downstream stages.
    \item \textbf{Normalization}: Images exceeding 4,096 pixels are downscaled with
          \texttt{INTER\_AREA} interpolation; images below 100 pixels are rejected.
    \item \textbf{Session management}: Each pipeline run generates a unique session ID
          (\texttt{sess\_YYYYMMDD\_HHMMSS\_\{uuid8\}}) and a config hash for
          reproducibility.
\end{itemize}

Stage~1 is marked as \emph{critical} -- the pipeline cannot proceed without valid input.
If ingestion fails, a \texttt{StageInputError} is raised rather than producing an
empty output.

% ============================================================
\subsection{Stage 2: Detection and Localization}
% ============================================================

Stage~2 detects chart regions in document images using a trained YOLOv8m model
\cite{jocher2023yolov8}. The detection pipeline operates as follows:

\begin{enumerate}
    \item Load each \texttt{CleanImage} and run YOLO inference (confidence threshold 0.5,
          NMS IoU threshold 0.45, input size $640 \times 640$).
    \item Filter detections by minimum area (100~pixels).
    \item For each valid detection, crop the chart region with 10-pixel padding to prevent
          text clipping, assign a unique \texttt{chart\_id}, and save the cropped image.
\end{enumerate}

The detector uses a \textbf{single-class strategy} (``chart'' vs.\ background) rather
than multi-class per chart type. This design decision is motivated by the observation
that YOLO excels at localization, while type classification benefits from a dedicated
architecture (ResNet-18, see \Cref{sec:system_s3}). The two-model cascade achieves
higher combined accuracy than a single multi-class detector.

Training uses conservative augmentation (no rotation or flipping, since charts are
orientation-sensitive) with mosaic, scale, HSV adjustments, and translation.

% ============================================================
\subsection{Stage 3: Structural Analysis (Extraction)}
\label{sec:system_s3}
% ============================================================

Stage~3 is the most complex component ($\sim$5,400 lines across 7 submodules). Its core
design principle is to treat charts as \textbf{collections of geometric entities} rather
than pixel patterns. \Cref{fig:s3_submodules} shows the submodule pipeline.

\begin{figure}[htbp]
\centering
% TikZ: Stage 3 Extraction Submodules
% Usage: % TikZ: Stage 3 Extraction Submodules
% Usage: % TikZ: Stage 3 Extraction Submodules
% Usage: \input{figures/tikz/tikz_stage3_submodules.tex}
% Requires: \usepackage{tikz} \usetikzlibrary{shapes.geometric, arrows.meta, positioning, fit, backgrounds}

\begin{tikzpicture}[
    node distance=0.4cm and 0.3cm,
    module/.style={
        rectangle, rounded corners=2pt, draw=black!60, thick,
        minimum width=2.6cm, minimum height=0.9cm,
        fill=#1!15, font=\scriptsize\bfseries, align=center
    },
    arrow/.style={-{Stealth[length=2.5mm]}, thick, draw=black!50},
    label/.style={font=\tiny\itshape, text=black!50, align=center},
    io/.style={
        rectangle, rounded corners=2pt, draw=black!40,
        minimum width=2.2cm, minimum height=0.6cm,
        fill=black!5, font=\scriptsize, align=center
    }
]

% Input
\node[io] (input) {Cropped Chart Image};

% Row 1: Preprocessing
\node[module=blue, below=0.5cm of input] (prep) {Preprocessor};
\node[label, right=0.05cm of prep] {Negative transform\\Adaptive threshold};

% Row 2: Skeleton + OCR (parallel)
\node[module=teal, below left=0.6cm and 0.3cm of prep] (skel) {Skeletonizer};
\node[module=orange, below right=0.6cm and 0.3cm of prep] (ocr) {OCR Engine};
\node[label, left=0.05cm of skel] {Lee algorithm\\Keypoint detection};
\node[label, right=0.05cm of ocr] {EasyOCR\\Role classification};

% Row 3: Vectorizer + Elements
\node[module=purple, below=0.6cm of skel] (vec) {Vectorizer};
\node[module=red, below=0.6cm of ocr] (elem) {Element Detector};
\node[label, left=0.05cm of vec] {Adaptive RDP\\Subpixel refinement};
\node[label, right=0.05cm of elem] {Bars, markers\\Pie slices};

% Row 4: Geometric Mapper + Classifier (merge)
\node[module=teal, below=0.8cm of vec, xshift=1.8cm] (geo) {Geometric Mapper};
\node[module=blue, right=0.4cm of geo] (cls) {ResNet-18\\Classifier};
\node[label, below=0.05cm of geo] {RANSAC calibration\\Pixel-to-value mapping};
\node[label, below=0.05cm of cls] {8 chart types\\94.14\% accuracy};

% Output
\node[io, below=0.8cm of geo, xshift=1.2cm] (output) {RawMetadata (Stage3Output)};

% Arrows
\draw[arrow] (input) -- (prep);
\draw[arrow] (prep) -| (skel);
\draw[arrow] (prep) -| (ocr);
\draw[arrow] (skel) -- (vec);
\draw[arrow] (ocr) -- (elem);
\draw[arrow] (vec) |- (geo);
\draw[arrow] (elem) |- (geo);
\draw[arrow] (ocr) |- (geo);
\draw[arrow] (prep) -| (cls);
\draw[arrow] (geo) -- (output);
\draw[arrow] (cls) |- (output);

\end{tikzpicture}

% Requires: \usepackage{tikz} \usetikzlibrary{shapes.geometric, arrows.meta, positioning, fit, backgrounds}

\begin{tikzpicture}[
    node distance=0.4cm and 0.3cm,
    module/.style={
        rectangle, rounded corners=2pt, draw=black!60, thick,
        minimum width=2.6cm, minimum height=0.9cm,
        fill=#1!15, font=\scriptsize\bfseries, align=center
    },
    arrow/.style={-{Stealth[length=2.5mm]}, thick, draw=black!50},
    label/.style={font=\tiny\itshape, text=black!50, align=center},
    io/.style={
        rectangle, rounded corners=2pt, draw=black!40,
        minimum width=2.2cm, minimum height=0.6cm,
        fill=black!5, font=\scriptsize, align=center
    }
]

% Input
\node[io] (input) {Cropped Chart Image};

% Row 1: Preprocessing
\node[module=blue, below=0.5cm of input] (prep) {Preprocessor};
\node[label, right=0.05cm of prep] {Negative transform\\Adaptive threshold};

% Row 2: Skeleton + OCR (parallel)
\node[module=teal, below left=0.6cm and 0.3cm of prep] (skel) {Skeletonizer};
\node[module=orange, below right=0.6cm and 0.3cm of prep] (ocr) {OCR Engine};
\node[label, left=0.05cm of skel] {Lee algorithm\\Keypoint detection};
\node[label, right=0.05cm of ocr] {EasyOCR\\Role classification};

% Row 3: Vectorizer + Elements
\node[module=purple, below=0.6cm of skel] (vec) {Vectorizer};
\node[module=red, below=0.6cm of ocr] (elem) {Element Detector};
\node[label, left=0.05cm of vec] {Adaptive RDP\\Subpixel refinement};
\node[label, right=0.05cm of elem] {Bars, markers\\Pie slices};

% Row 4: Geometric Mapper + Classifier (merge)
\node[module=teal, below=0.8cm of vec, xshift=1.8cm] (geo) {Geometric Mapper};
\node[module=blue, right=0.4cm of geo] (cls) {ResNet-18\\Classifier};
\node[label, below=0.05cm of geo] {RANSAC calibration\\Pixel-to-value mapping};
\node[label, below=0.05cm of cls] {8 chart types\\94.14\% accuracy};

% Output
\node[io, below=0.8cm of geo, xshift=1.2cm] (output) {RawMetadata (Stage3Output)};

% Arrows
\draw[arrow] (input) -- (prep);
\draw[arrow] (prep) -| (skel);
\draw[arrow] (prep) -| (ocr);
\draw[arrow] (skel) -- (vec);
\draw[arrow] (ocr) -- (elem);
\draw[arrow] (vec) |- (geo);
\draw[arrow] (elem) |- (geo);
\draw[arrow] (ocr) |- (geo);
\draw[arrow] (prep) -| (cls);
\draw[arrow] (geo) -- (output);
\draw[arrow] (cls) |- (output);

\end{tikzpicture}

% Requires: \usepackage{tikz} \usetikzlibrary{shapes.geometric, arrows.meta, positioning, fit, backgrounds}

\begin{tikzpicture}[
    node distance=0.4cm and 0.3cm,
    module/.style={
        rectangle, rounded corners=2pt, draw=black!60, thick,
        minimum width=2.6cm, minimum height=0.9cm,
        fill=#1!15, font=\scriptsize\bfseries, align=center
    },
    arrow/.style={-{Stealth[length=2.5mm]}, thick, draw=black!50},
    label/.style={font=\tiny\itshape, text=black!50, align=center},
    io/.style={
        rectangle, rounded corners=2pt, draw=black!40,
        minimum width=2.2cm, minimum height=0.6cm,
        fill=black!5, font=\scriptsize, align=center
    }
]

% Input
\node[io] (input) {Cropped Chart Image};

% Row 1: Preprocessing
\node[module=blue, below=0.5cm of input] (prep) {Preprocessor};
\node[label, right=0.05cm of prep] {Negative transform\\Adaptive threshold};

% Row 2: Skeleton + OCR (parallel)
\node[module=teal, below left=0.6cm and 0.3cm of prep] (skel) {Skeletonizer};
\node[module=orange, below right=0.6cm and 0.3cm of prep] (ocr) {OCR Engine};
\node[label, left=0.05cm of skel] {Lee algorithm\\Keypoint detection};
\node[label, right=0.05cm of ocr] {EasyOCR\\Role classification};

% Row 3: Vectorizer + Elements
\node[module=purple, below=0.6cm of skel] (vec) {Vectorizer};
\node[module=red, below=0.6cm of ocr] (elem) {Element Detector};
\node[label, left=0.05cm of vec] {Adaptive RDP\\Subpixel refinement};
\node[label, right=0.05cm of elem] {Bars, markers\\Pie slices};

% Row 4: Geometric Mapper + Classifier (merge)
\node[module=teal, below=0.8cm of vec, xshift=1.8cm] (geo) {Geometric Mapper};
\node[module=blue, right=0.4cm of geo] (cls) {ResNet-18\\Classifier};
\node[label, below=0.05cm of geo] {RANSAC calibration\\Pixel-to-value mapping};
\node[label, below=0.05cm of cls] {8 chart types\\94.14\% accuracy};

% Output
\node[io, below=0.8cm of geo, xshift=1.2cm] (output) {RawMetadata (Stage3Output)};

% Arrows
\draw[arrow] (input) -- (prep);
\draw[arrow] (prep) -| (skel);
\draw[arrow] (prep) -| (ocr);
\draw[arrow] (skel) -- (vec);
\draw[arrow] (ocr) -- (elem);
\draw[arrow] (vec) |- (geo);
\draw[arrow] (elem) |- (geo);
\draw[arrow] (ocr) |- (geo);
\draw[arrow] (prep) -| (cls);
\draw[arrow] (geo) -- (output);
\draw[arrow] (cls) |- (output);

\end{tikzpicture}

\caption{Stage~3 extraction pipeline: 7 submodules process the cropped chart image
    sequentially. Each submodule has an independent configuration and can be toggled on/off.}
\label{fig:s3_submodules}
\end{figure}

\subsubsection{Preprocessor}

The preprocessor converts chart images to optimized binary masks. A key technique is the
\textbf{negative image transform}: $I_{\text{neg}} = 255 - I$. This inversion aligns
chart strokes (thin, dark lines) with morphological operations that expect bright
foreground, improving skeleton quality by approximately 30\%. Adaptive Gaussian
thresholding ($T(x,y) = \mu_G(x,y) - C$, block size 11, $C=2$) handles non-uniform
lighting in scanned documents.

\subsubsection{Skeletonizer}

The Lee topology-preserving thinning algorithm \cite{lee1994skeleton} reduces thick
strokes to 1-pixel width while maintaining 8-connectivity and homotopy (no holes or
junctions are created or destroyed). Gap filling via morphological closing
(\texttt{kernel=3$\times$3}) precedes skeletonization to repair small breaks in chart
lines. Junction detection classifies skeleton pixels as endpoints (1~neighbor),
junctions (3+ neighbors), or intermediate (2~neighbors).

\subsubsection{Vectorizer}

The Ramer-Douglas-Peucker (RDP) algorithm \cite{douglas1973rdp, ramer1972rdp}
recursively simplifies polylines by removing points within tolerance $\epsilon$.
Geo-SLM uses an \textbf{adaptive epsilon}:
\begin{equation}
    \epsilon = \epsilon_{\text{base}} \cdot \sqrt{\frac{L}{L_{\text{ref}}}}
\end{equation}
where $L$ is the polyline length and $L_{\text{ref}}$ is a reference length. This
prevents over-simplification of long curves and under-simplification of short ones.
Sub-pixel refinement uses weighted centroids for fractional pixel positions.

\subsubsection{OCR Engine}

PaddleOCR \cite{du2020paddleocr} (with EasyOCR \cite{jaided2020easyocr} as fallback)
extracts all text, which is then assigned spatial roles based on position: title
(top 15\%), X-axis label (bottom 15\%), Y-axis label (left 15\%), legend
(top-right corner), and data values (near elements). Content-aware keyword matching
supplements positional heuristics.

\subsubsection{Element Detector}

Three detection methods are combined for bar detection: watershed segmentation for
touching bars, projection profile analysis for alignment, and morphological connected
components. Scatter plot markers use Hough circle transform. Pie slices use contour
analysis with arc fitting. Color clustering (K-means, $k=8$) identifies dominant
colors per element.

\subsubsection{Geometric Mapper}

The geometric mapper calibrates axes and converts pixel coordinates to data values.
The calibration pipeline:

\begin{enumerate}
    \item Match OCR tick labels to detected tick positions along axes.
    \item Build (pixel, value) correspondence pairs.
    \item Fit a calibration model using RANSAC \cite{fischler1981ransac} (default),
          which tolerates up to 40\% outlier rejection from OCR errors.
    \item Compute $R^2$ as calibration confidence.
\end{enumerate}

The linear mapping function (with Y-axis inversion for image coordinates):
\begin{equation}
    v_y = y_{\max} - \frac{p_y - p_{\text{top}}}{p_{\text{bottom}} - p_{\text{top}}}
          \cdot (y_{\max} - y_{\min})
\end{equation}
Logarithmic scale is detected automatically when tick values follow a geometric
progression.

\subsubsection{Classifier}

A cascade of three classifiers determines chart type:

\begin{enumerate}
    \item \textbf{ResNet-18} \cite{he2016resnet}: CNN on $256 \times 256$ grayscale input,
          8 classes, 94.14\% test accuracy, 6.90ms ONNX inference (CPU).
    \item \textbf{ML Classifier}: Random Forest on structural features ($\sim$70\% accuracy).
    \item \textbf{Rule-based}: Heuristics (has bars? has slices?) as final fallback
          ($\sim$50\% accuracy).
\end{enumerate}

The overall extraction confidence is a weighted average:
\begin{equation}
    C_{\text{overall}} = 0.3 \cdot C_{\text{class}} + 0.25 \cdot C_{\text{ocr}}
                       + 0.25 \cdot C_{\text{axis}} + 0.2 \cdot C_{\text{elem}}
\end{equation}

% ============================================================
\subsection{Stage 4: Semantic Reasoning}
% ============================================================

Stage~4 refines raw extracted metadata using AI reasoning. The processing pipeline
per chart consists of four steps:

\textbf{Step 1: Geometric Value Mapping.}
The \texttt{GeometricValueMapper} (764 lines) applies the calibration from Stage~3 to
convert all pixel coordinates to data values. Confidence filtering ensures mapping is
applied only when calibration confidence $\geq 0.3$. Value clamping restricts results
to the detected axis range to prevent extrapolation errors.

\textbf{Step 2: Prompt Construction.}
The \texttt{GeminiPromptBuilder} (833 lines) constructs structured prompts with explicit
sections: \textsc{Chart\_Metadata}, \textsc{OCR\_Texts}, \textsc{Elements},
\textsc{Axis\_Info}, \textsc{Mapped\_Values}, and \textsc{Instructions}. Anti-hallucination
rules are embedded: ``Only report values that appear in OCR\_TEXTS or MAPPED\_VALUES;
if uncertain, report null rather than guess.''

\textbf{Step 3: AI Routing.}
The \texttt{AIRouterEngine} delegates to the AI Router (see \Cref{sec:ai_routing}).
The router walks the task-specific fallback chain until a provider returns a response
with confidence above threshold (0.7).

\textbf{Step 4: Post-processing.}
The AI response is parsed into structured \texttt{RefinedChartData}: title, axis labels,
named data series with $(x, y)$ points, natural language description, and confidence
scores.

A pure rule-based fallback is available when all AI providers fail, applying pattern-based
OCR corrections and direct use of geometric mapping values without LLM refinement.

% ============================================================
\subsection{Stage 5: Insight and Reporting}
% ============================================================

Stage~5 generates human-readable insights from structured chart data and produces
final output files. Four insight types are supported:

\begin{itemize}
    \item \textbf{Trend detection}: Linear regression slope classification
          (increasing, decreasing, stable, volatile), with $R^2$ as confidence.
    \item \textbf{Comparison}: Cross-series analysis (dominant series, largest gap,
          crossover points) for charts with $\geq 2$ data series.
    \item \textbf{Anomaly detection}: Z-score method ($|z| > 2.0$) applied per-series.
    \item \textbf{Summary}: Template-based text combining chart type, axis labels,
          data range, and trend information.
\end{itemize}

Output is produced in two formats: (1) a structured JSON file containing the full
\texttt{PipelineResult} with session metadata, all chart results, insights, and
source traceability; and (2) a Markdown report with per-chart sections. Failed charts
are included with error insights rather than dropped, preserving pipeline completeness.

% ============================================================
\subsection{AI Routing Layer}
\label{sec:ai_routing}
% ============================================================

\Cref{fig:ai_router} illustrates the AI routing architecture.

\begin{figure}[htbp]
\centering
% TikZ: AI Router with Fallback Chains
% Usage: % TikZ: AI Router with Fallback Chains
% Usage: % TikZ: AI Router with Fallback Chains
% Usage: \input{figures/tikz/tikz_ai_router.tex}
% Requires: \usepackage{tikz} \usetikzlibrary{shapes.geometric, arrows.meta, positioning, fit, backgrounds, calc}

\begin{tikzpicture}[
    node distance=0.6cm and 0.8cm,
    block/.style={
        rectangle, rounded corners=3pt, draw=black!60, thick,
        minimum width=2.4cm, minimum height=0.9cm,
        font=\footnotesize\bfseries, align=center
    },
    adapter/.style={
        rectangle, rounded corners=3pt, draw=#1!60!black, thick,
        minimum width=2.2cm, minimum height=0.8cm,
        fill=#1!15, font=\footnotesize, align=center
    },
    task/.style={
        rectangle, rounded corners=2pt, draw=black!40,
        minimum width=2.0cm, minimum height=0.6cm,
        fill=yellow!10, font=\scriptsize, align=center
    },
    arrow/.style={-{Stealth[length=2.5mm]}, thick, draw=black!50},
    dasharrow/.style={-{Stealth[length=2.5mm]}, dashed, draw=red!50, thick},
    label/.style={font=\scriptsize\itshape, text=black!50}
]

% Stage 4 caller
\node[block, fill=purple!15] (caller) {Stage 4\\Reasoning};

% Router
\node[block, fill=orange!20, right=1.5cm of caller] (router) {AI Router\\resolve()};

% Adapters
\node[adapter=teal, right=1.5cm of router, yshift=1.0cm] (slm) {Local SLM\\Adapter};
\node[adapter=blue, right=1.5cm of router] (gemini) {Gemini\\Adapter};
\node[adapter=red, right=1.5cm of router, yshift=-1.0cm] (openai) {OpenAI\\Adapter};

% Standardized output
\node[block, fill=green!15, right=1.2cm of gemini] (result) {Reasoning\\Result};

% Arrows: caller -> router
\draw[arrow] (caller) -- node[above, label] {TaskType} (router);

% Router -> adapters
\draw[arrow] (router) -- (slm);
\draw[arrow] (router) -- (gemini);
\draw[arrow] (router) -- (openai);

% Adapters -> result
\draw[arrow] (slm) -- (result);
\draw[arrow] (gemini) -- (result);
\draw[arrow] (openai) -- (result);

% Fallback arrows
\draw[dasharrow] (slm.south) -- node[right, label, text=red!60] {fallback} (gemini.north);
\draw[dasharrow] (gemini.south) -- node[right, label, text=red!60] {fallback} (openai.north);

% Task type examples below router
\node[task, below=0.8cm of router, xshift=-1.2cm] (t1) {OCR Correction};
\node[task, right=0.2cm of t1] (t2) {Value Mapping};
\node[task, below=0.3cm of t1, xshift=0.6cm] (t3) {Chart Reasoning};

\node[label, above=0.05cm of t1.north west, anchor=south west] {Task Types:};

% Priority labels
\node[label, left=0.1cm of slm, text=teal!70] {Priority 1};
\node[label, left=0.1cm of gemini, text=blue!70] {Priority 2};
\node[label, left=0.1cm of openai, text=red!70] {Priority 3};

\end{tikzpicture}

% Requires: \usepackage{tikz} \usetikzlibrary{shapes.geometric, arrows.meta, positioning, fit, backgrounds, calc}

\begin{tikzpicture}[
    node distance=0.6cm and 0.8cm,
    block/.style={
        rectangle, rounded corners=3pt, draw=black!60, thick,
        minimum width=2.4cm, minimum height=0.9cm,
        font=\footnotesize\bfseries, align=center
    },
    adapter/.style={
        rectangle, rounded corners=3pt, draw=#1!60!black, thick,
        minimum width=2.2cm, minimum height=0.8cm,
        fill=#1!15, font=\footnotesize, align=center
    },
    task/.style={
        rectangle, rounded corners=2pt, draw=black!40,
        minimum width=2.0cm, minimum height=0.6cm,
        fill=yellow!10, font=\scriptsize, align=center
    },
    arrow/.style={-{Stealth[length=2.5mm]}, thick, draw=black!50},
    dasharrow/.style={-{Stealth[length=2.5mm]}, dashed, draw=red!50, thick},
    label/.style={font=\scriptsize\itshape, text=black!50}
]

% Stage 4 caller
\node[block, fill=purple!15] (caller) {Stage 4\\Reasoning};

% Router
\node[block, fill=orange!20, right=1.5cm of caller] (router) {AI Router\\resolve()};

% Adapters
\node[adapter=teal, right=1.5cm of router, yshift=1.0cm] (slm) {Local SLM\\Adapter};
\node[adapter=blue, right=1.5cm of router] (gemini) {Gemini\\Adapter};
\node[adapter=red, right=1.5cm of router, yshift=-1.0cm] (openai) {OpenAI\\Adapter};

% Standardized output
\node[block, fill=green!15, right=1.2cm of gemini] (result) {Reasoning\\Result};

% Arrows: caller -> router
\draw[arrow] (caller) -- node[above, label] {TaskType} (router);

% Router -> adapters
\draw[arrow] (router) -- (slm);
\draw[arrow] (router) -- (gemini);
\draw[arrow] (router) -- (openai);

% Adapters -> result
\draw[arrow] (slm) -- (result);
\draw[arrow] (gemini) -- (result);
\draw[arrow] (openai) -- (result);

% Fallback arrows
\draw[dasharrow] (slm.south) -- node[right, label, text=red!60] {fallback} (gemini.north);
\draw[dasharrow] (gemini.south) -- node[right, label, text=red!60] {fallback} (openai.north);

% Task type examples below router
\node[task, below=0.8cm of router, xshift=-1.2cm] (t1) {OCR Correction};
\node[task, right=0.2cm of t1] (t2) {Value Mapping};
\node[task, below=0.3cm of t1, xshift=0.6cm] (t3) {Chart Reasoning};

\node[label, above=0.05cm of t1.north west, anchor=south west] {Task Types:};

% Priority labels
\node[label, left=0.1cm of slm, text=teal!70] {Priority 1};
\node[label, left=0.1cm of gemini, text=blue!70] {Priority 2};
\node[label, left=0.1cm of openai, text=red!70] {Priority 3};

\end{tikzpicture}

% Requires: \usepackage{tikz} \usetikzlibrary{shapes.geometric, arrows.meta, positioning, fit, backgrounds, calc}

\begin{tikzpicture}[
    node distance=0.6cm and 0.8cm,
    block/.style={
        rectangle, rounded corners=3pt, draw=black!60, thick,
        minimum width=2.4cm, minimum height=0.9cm,
        font=\footnotesize\bfseries, align=center
    },
    adapter/.style={
        rectangle, rounded corners=3pt, draw=#1!60!black, thick,
        minimum width=2.2cm, minimum height=0.8cm,
        fill=#1!15, font=\footnotesize, align=center
    },
    task/.style={
        rectangle, rounded corners=2pt, draw=black!40,
        minimum width=2.0cm, minimum height=0.6cm,
        fill=yellow!10, font=\scriptsize, align=center
    },
    arrow/.style={-{Stealth[length=2.5mm]}, thick, draw=black!50},
    dasharrow/.style={-{Stealth[length=2.5mm]}, dashed, draw=red!50, thick},
    label/.style={font=\scriptsize\itshape, text=black!50}
]

% Stage 4 caller
\node[block, fill=purple!15] (caller) {Stage 4\\Reasoning};

% Router
\node[block, fill=orange!20, right=1.5cm of caller] (router) {AI Router\\resolve()};

% Adapters
\node[adapter=teal, right=1.5cm of router, yshift=1.0cm] (slm) {Local SLM\\Adapter};
\node[adapter=blue, right=1.5cm of router] (gemini) {Gemini\\Adapter};
\node[adapter=red, right=1.5cm of router, yshift=-1.0cm] (openai) {OpenAI\\Adapter};

% Standardized output
\node[block, fill=green!15, right=1.2cm of gemini] (result) {Reasoning\\Result};

% Arrows: caller -> router
\draw[arrow] (caller) -- node[above, label] {TaskType} (router);

% Router -> adapters
\draw[arrow] (router) -- (slm);
\draw[arrow] (router) -- (gemini);
\draw[arrow] (router) -- (openai);

% Adapters -> result
\draw[arrow] (slm) -- (result);
\draw[arrow] (gemini) -- (result);
\draw[arrow] (openai) -- (result);

% Fallback arrows
\draw[dasharrow] (slm.south) -- node[right, label, text=red!60] {fallback} (gemini.north);
\draw[dasharrow] (gemini.south) -- node[right, label, text=red!60] {fallback} (openai.north);

% Task type examples below router
\node[task, below=0.8cm of router, xshift=-1.2cm] (t1) {OCR Correction};
\node[task, right=0.2cm of t1] (t2) {Value Mapping};
\node[task, below=0.3cm of t1, xshift=0.6cm] (t3) {Chart Reasoning};

\node[label, above=0.05cm of t1.north west, anchor=south west] {Task Types:};

% Priority labels
\node[label, left=0.1cm of slm, text=teal!70] {Priority 1};
\node[label, left=0.1cm of gemini, text=blue!70] {Priority 2};
\node[label, left=0.1cm of openai, text=red!70] {Priority 3};

\end{tikzpicture}

\caption{AI Router with Adapter Pattern and per-task fallback chains.
    Each provider implements the common \texttt{BaseAIAdapter} interface.}
\label{fig:ai_router}
\end{figure}

All AI providers implement \texttt{BaseAIAdapter}, an abstract base class defining
\texttt{reason()}, \texttt{correct\_ocr()}, and \texttt{health\_check()} methods.
Concrete adapters (\texttt{GeminiAdapter}, \texttt{OpenAIAdapter},
\texttt{LocalSLMAdapter}) handle provider-specific API calls, authentication, and
error wrapping.

\Cref{tab:ai_providers} shows the configured providers and their roles.

\begin{table}[htbp]
\centering
\caption{AI provider configuration.}
\label{tab:ai_providers}
\begin{tabular}{@{}lllll@{}}
\toprule
\textbf{Provider} & \textbf{Model} & \textbf{Use Case} & \textbf{Priority} & \textbf{Cost} \\
\midrule
Local SLM & Qwen-2.5-1.5B    & Default reasoning   & Primary    & Free \\
Gemini    & gemini-2.0-flash  & Complex charts       & Fallback 1 & API \\
OpenAI    & gpt-4o-mini       & Alternative          & Fallback 2 & API \\
\bottomrule
\end{tabular}
\end{table}

The \texttt{AIRouter} resolves a (adapter, model\_id) pair for each task type by:
(1)~checking configuration overrides, (2)~walking the fallback chain for the task type,
(3)~verifying API key availability and provider health, and (4)~returning the first
healthy provider. A \textbf{local-only mode} restricts routing to the local SLM for
privacy-sensitive or offline deployments.

% ============================================================
\subsection{Error Handling and Extensibility}
% ============================================================

The system uses a hierarchical exception system rooted at \texttt{ChartAnalysisError}.
Each exception carries \texttt{stage}, \texttt{recoverable}, and
\texttt{fallback\_available} flags, enabling the orchestrator to decide whether to
continue, retry, or abort.

The modular architecture provides five extension points:
\begin{itemize}
    \item \textbf{New chart type}: Add an enum value to \texttt{ChartType} and element
          detection rules.
    \item \textbf{New AI provider}: Implement \texttt{BaseAIAdapter} and register in the
          router configuration.
    \item \textbf{New pipeline stage}: Implement \texttt{BaseStage[In, Out]} and register
          in the orchestrator.
    \item \textbf{New output format}: Add to \texttt{OutputFormat} enum and write a
          formatter in Stage~5.
    \item \textbf{New language}: Change OCR config, add keyword dictionaries, create
          prompt template variants, and fine-tune SLM with target-language data --
          all without modifying core pipeline logic.
    \item \textbf{Configuration override}: Add a YAML key; OmegaConf handles hierarchical
          merging.
\end{itemize}

% ============================================================
\subsection{Vietnamese Language Implementation}
\label{sec:viet_impl}
% ============================================================

This subsection details the concrete implementation path for Vietnamese chart
understanding, demonstrating how the extension module architecture supports language
localization at each pipeline stage.

\subsubsection{Language-Agnostic vs.\ Language-Dependent Components}

\Cref{tab:lang_analysis} classifies each pipeline component by its language dependency.
The classification reveals that the majority of the system -- the entire geometric
analysis core -- is inherently language-agnostic.

\begin{table}[htbp]
\centering
\caption{Language dependency analysis of pipeline components.}
\label{tab:lang_analysis}
\begin{tabular}{@{}lllp{5.5cm}@{}}
\toprule
\textbf{Component} & \textbf{Stage} & \textbf{Language Dep.} & \textbf{Rationale} \\
\midrule
PDF Renderer           & S1 & Agnostic  & Pixel-level rendering, no text parsing \\
Blur Detector          & S1 & Agnostic  & Laplacian variance is language-free \\
YOLO Detector          & S2 & Agnostic  & Visual bounding box detection \\
Preprocessor           & S3 & Agnostic  & Negative transform, thresholding \\
Skeletonizer           & S3 & Agnostic  & Topology-preserving thinning \\
Vectorizer (RDP)       & S3 & Agnostic  & Polyline simplification \\
Element Detector       & S3 & Agnostic  & Watershed, Hough, contour analysis \\
Geometric Mapper       & S3 & Agnostic  & Pixel-to-value calibration (numeric) \\
ResNet-18 Classifier   & S3 & Agnostic  & Visual feature classification \\
\midrule
OCR Engine             & S3 & \textbf{Config}    & Language set via \texttt{config/models.yaml} \\
OCR Role Classifier    & S3 & \textbf{Keywords}  & English keyword lists, needs Vietnamese extension \\
Prompt Builder         & S4 & \textbf{Templates} & English prompts, needs Vietnamese variants \\
AI Adapters            & S4 & Agnostic  & Pass-through interface, language in prompts \\
AI Router              & S4 & Agnostic  & Provider selection, language-transparent \\
Insight Generator      & S5 & \textbf{Templates} & English summary templates \\
QA Generator           & Data & \textbf{Prompts}   & English generation instructions \\
\bottomrule
\end{tabular}
\end{table}

Out of 17 major components, only 4 require modification (OCR role classifier, prompt
builder, insight generator, QA generator) and 1 requires configuration change (OCR
engine). The remaining 12 components, including all geometric analysis modules, are
fully language-agnostic.

\subsubsection{OCR Engine: Vietnamese Configuration Path}

The \texttt{OCRConfig} class defines a \texttt{languages} parameter (default:
\texttt{["en"]}). Vietnamese support is activated through configuration without code
changes:

\begin{verbatim}
# config/models.yaml -- Vietnamese configuration
ocr:
  paddleocr:
    lang: "vi"          # Vietnamese recognition weights
  easyocr:
    languages: ["vi"]   # Multi-language support
  tesseract:
    lang: "vie"         # Tesseract Vietnamese tessdata
\end{verbatim}

The OCR engine initialization code (\texttt{ocr\_engine.py}) already supports this:
PaddleOCR passes \texttt{lang=self.config.languages[0]}, and EasyOCR passes the full
language list to \texttt{easyocr.Reader()}. For bilingual charts (mixed English and
Vietnamese), the configuration \texttt{["en", "vi"]} enables simultaneous recognition.

\subsubsection{Stage 3: Bilingual Role Classification}

The \texttt{\_classify\_role()} method in the OCR engine uses keyword dictionaries for
semantic role assignment. The Vietnamese extension adds parallel dictionaries:

\begin{verbatim}
# Vietnamese axis label keywords (additive extension)
axis_label_keywords_vi = [
    'so luong', 'ty le', 'gia tri', 'doanh thu',
    'thoi gian', 'nam', 'thang', 'ngay',
    'gia', 'chi phi', 'loi nhuan', 'diem so',
    'do chinh xac', 'sai so', 'mat mat',
]

# Vietnamese legend keywords
legend_keywords_vi = [
    'nhom', 'loai', 'phan loai', 'chuoi',
    'huan luyen', 'kiem tra', 'du doan',
    'mo hinh', 'phuong phap', 'thuat toan',
]
\end{verbatim}

The position-based fallback logic (assigning roles based on text position within the
chart image -- top 15\% for titles, bottom 15\% for X-axis, etc.) remains effective
regardless of language, providing a robust baseline when keyword matching fails.

\subsubsection{Stage 4: Vietnamese Prompt Architecture}

The AI adapter interface (\texttt{BaseAIAdapter.reason()}) accepts
\texttt{system\_prompt} and \texttt{user\_prompt} as string parameters, making it
inherently language-transparent. Vietnamese support is implemented entirely at the
prompt layer:

\begin{enumerate}
    \item \textbf{Vietnamese prompt template files}: New files
          \texttt{prompts/reasoning\_vi.txt}, \texttt{description\_vi.txt}, and
          \texttt{ocr\_correction\_vi.txt} contain Vietnamese system instructions with
          Vietnamese chart analysis terminology.
    \item \textbf{Language routing}: The \texttt{PromptConfig.output\_language} field
          (already defined as \texttt{"english"} by default) is wired into the prompt
          selection logic to choose between English and Vietnamese template sets.
    \item \textbf{SLM compatibility}: The Qwen-2.5 model family supports Vietnamese
          natively in its tokenizer, ensuring efficient token usage for Vietnamese text
          generation without excessive subword splitting.
\end{enumerate}

The anti-hallucination constraints (``only report values from OCR\_TEXTS or
MAPPED\_VALUES'') are preserved in Vietnamese prompts, translated to equivalent
Vietnamese instructions to maintain output grounding.

\subsubsection{Data Factory: Vietnamese QA Generation}

The QA generation pipeline in \texttt{tools/data\_factory/} produces training data by
sending chart features to Gemini API with generation instructions. Vietnamese QA
generation requires:

\begin{enumerate}
    \item Rewriting the system prompt in \texttt{qa\_generator.py} to enforce Vietnamese
          output: all questions, answers, and explanations must be in Vietnamese.
    \item Expanding the question taxonomy with Vietnamese-specific terminology covering
          chart analysis concepts (e.g.,
          ``xu hướng tăng/giảm'',
          ``giá trị lớn nhất'',
          ``so sánh giữa các nhóm'').
    \item Targeting 5,000--10,000 Vietnamese QA samples to supplement the existing
          268,799 English samples.
    \item Validating generated data through the existing verification pipeline with
          Vietnamese-aware quality checks.
\end{enumerate}

The combined English + Vietnamese training data enables the SLM to handle both languages,
with the English foundation providing geometric reasoning knowledge and the Vietnamese
data providing language-specific output capability.